\section{Discussion}
\label{sec:discussion}

\subsection{Key Findings and Their Interpretation}

\textbf{Finding 1: The structural efficiency framework is validated end-to-end.}
The most secure contribution of this paper is the measurement framework itself.
All components execute correctly, produce interpretable metrics, and handle the
full statistical pipeline. This means researchers can immediately apply the
framework to their own GRPO/DPO checkpoint archives --- the infrastructure cost
is already paid.

\textbf{Finding 2: The raw/proportion dissociation challenges the selective-reallocation mechanism.}
The most intellectually provocative preliminary result is the divergence between
GRPO's raw semantic edit distance advantage (+250\% on PoC data; h-e1 synthetic
PoC data, not real training output) and its near-zero SEP advantage ($\approx$0.237 for
both methods; h-m1 preliminary, $n_\text{eff}\approx2$). These two metrics ask different
questions: raw edit distance asks ``how much does the policy change structurally
in absolute terms?''; SEP asks ``what fraction of all changes are structural?''

Three competing interpretations deserve investigation:

\textbf{(a) The mechanism claim (P2) is false even if the existence claim (P1) holds.}
GRPO may produce more aggressive code restructuring (more total edits, including
more semantic edits in absolute terms) without proportionally concentrating edits
on semantic nodes. This would mean execution reward increases structural activity
broadly, not selectively.

\textbf{(b) Checkpoint aliasing masks a real SEP difference.}
With $n_\text{eff} \approx 2$, the Mann-Whitney test has essentially no statistical
power. A corrected run with 10+ diverse checkpoints might reveal $\text{SEP}_\text{GRPO} > \text{SEP}_\text{DPO}$
with adequate power. This interpretation cannot be ruled out
and motivates the corrected experimental protocol (Section~\ref{sec:corrected_protocol}).

\textbf{(c) Scale-level equilibration of SEP at 7B.}
At 7B parameter scale with the specific KL budget used, both methods may converge
to similar proportional distributions over node types, even when absolute edit
distances differ. This would be a scale-specific finding that motivates testing
at different model sizes (1.3B, 13B).

Distinguishing these interpretations requires a corrected run with real training
($\geq$10 GRPO checkpoints from $\geq$1000 training steps) and execution-oracle DPO pairs.
The framework is ready; the experiment needs to be run.

\textbf{Finding 3: Checkpoint aliasing is a confound not, to our knowledge, previously documented in the RL fine-tuning literature.}
This finding has value independent of any hypothesis about GRPO vs.\ DPO. Any
researcher running checkpoint-comparison analysis in RL fine-tuning --- for code,
math, or general NLP --- should add checkpoint diversity pre-flight checks to their
experimental pipeline.

\subsection{Limitations}
\label{sec:limitations}

\textbf{L1: Synthetic data in h-e1 gate evaluation.}
The proof-of-concept (h-e1) used hand-crafted code completions with guaranteed
GRPO structural advantage. The bootstrap CI [4.65, 8.73] and mean differential 6.50
are valid measurements of the infrastructure's detection capability, but cannot
be cited as empirical evidence that real GRPO training produces higher structural efficiency than DPO.

\textbf{L2: Underpowered SEP analysis ($n_\text{effective} \approx 2$).}
The h-m1 SEP comparison is entirely underpowered. The Mann-Whitney test
($p = 0.4248$) and the near-zero effect size ($-0.0072$) cannot be interpreted as
evidence for or against the mechanism hypothesis. We report these numbers for
transparency and document the aliasing root cause; they should not be cited as
empirical findings about GRPO vs.\ DPO behavior.

\textbf{L3: Single base model and single programming language.}
All experiments use DeepSeek-Coder-7B-instruct-v1.5 on Python (HumanEval+/MBPP+).
The structural efficiency framework is designed to be model-agnostic and
language-agnostic, but empirical generalization has not been tested.

\textbf{L4: DPO preference pairs not execution-oracle labeled.}
The h-e1 DPO implementation used stub preference pairs (\texttt{return None} as
rejected completion) rather than genuine model-generated failed solutions
labeled by evalplus. Real execution-oracle DPO pairs require sampling model
outputs and labeling via execution.

\textbf{L5: SEP not validated against functional outcomes.}
The structural efficiency metric and SEP have not been validated against functional
correctness measures (pass@1, ECE, OOD transfer). The framework measures \emph{structural activity}
of policy movement, not the \emph{correctness} of that movement. This validation is part of the future work
outlined in Section~\ref{sec:future}.

\subsection{Corrected Experimental Protocol}
\label{sec:corrected_protocol}

For any team seeking to reproduce or extend this work with valid empirical results:

\begin{enumerate}
  \item \textbf{Enforce checkpoint diversity:} \texttt{save\_steps=100} for a 1000-step run yields
    10 GRPO checkpoints (steps 100--1000). Assert \texttt{len(unique\_checkpoints) >= 10}
    before h-m1 analysis.
  \item \textbf{Fix mock data:} Replace hard-coded GRPO code completions in \texttt{smoke\_test\_experiment.py}
    with real GRPOTrainer outputs evaluated by evalplus.
  \item \textbf{Fix DPO pairs:} Implement \texttt{generate\_dpo\_pairs()} to sample model outputs
    and label via evalplus execution (passing = preferred, failing = rejected).
  \item \textbf{Use KL tolerance = 0.05} ($\pm 5\%$) as specified, not 0.15. With 10+ checkpoints,
    sufficient pairs will be found even at stricter tolerance.
  \item \textbf{Reduce scope:} Test the H-E1 $\rightarrow$ H-M1 chain before attempting H-M2--H-M4.
    Confirm real-training SEP direction before proceeding to mediation analysis.
\end{enumerate}

\subsection{Broader Impact}

\textbf{Positive impacts:} The structural efficiency framework enables more principled
comparison of post-training methods --- researchers can now ask not only ``which
method scores higher?'' but ``which method induces richer structural learning per
unit of KL cost?'' The checkpoint aliasing finding will prevent silent data corruption
in future checkpoint-comparison studies across the RL fine-tuning literature.

\textbf{Potential risks and mitigations:} The structural efficiency metric could be
gamed if used as a direct training reward. We recommend the metric be used exclusively
for post-hoc analysis, not as a training signal, until its relationship to
functional correctness and out-of-distribution generalization is empirically
established.
